\clearpage\SourceLeaf{526}\pdfbookmark[1]{Für Ihr gütiges Schreiben}{gauss-dirichlet-fur-ihr-gutiges-schreiben}
Für Ihr gütiges Schreiben, und die gefällige Uebersendung Ihrer beiden Abhandlungen statte ich Ihnen, mein hochgeschätzter Freund, meinen verbindlichsten Dank ab. Ich sehe mit Vergnügen das steigende Interesse, welches man gegenwärtig an den Untersuchungen der Höhern Arithmetik zu nehmen anfängt. Die glückliche Art, wie Sie das zweite auf die biquadratische Residualität der Zahl 2 aus dem ersten ableiten, hat mir sehr wohl gefallen.
\par Vermuthlich hat jetzt der 6. Band unsrer Commentationen seinen Weg nach Breslau gefunden, und meine Commentatio prima über die biquadratischen Reste wird Ihnen also wol gegenwärtig bekannt sein: wenn sich eine Gelegenheit darbieten sollte, würde ich auch mit Vergnügen Ihnen einen besondern Abdruck derselben übersenden. Ich hätte unter mehrern Beweisarten für das darin vorkommende Theorem wählen können; es wird Ihnen aber nicht entgehen, warum ich den daselbst ausgeführten hier vorgezogen habe, hauptsächlich nemlich, weil die Classification von 2 bei denjenigen Moduln, für welche es quadratischer Nichtrest ist (unter $B$ oder $D$) als ein wesentlicher integrirender Theil des Theorems betrachtet werden muss, auf welchen die meisten andern Beweisarten nicht anwendbar scheinen.
\par Die ganze Untersuchung, deren Stoff ich schon seit 23 Jahren vollständig besitze, die Beweise der Haupttheoreme aber (zu welchen das in der ersten Commentation noch \textit{nicht} zu rechnen ist) seit etwa 14 Jahren \textemdash{} (obwol ich wünsche und hoffe, an letztern, den Beweisen, noch einiges vereinfachen zu können) \textemdash{} habe ich auf ungefähr 3 Abhandlungen berechnet. Mit der Abfassung der zweiten habe ich bereits jetzt einen Anfang gemacht, und hoffe sie in nicht langer Zeit zu vollenden, falls nicht die neuerdings mir wieder aufgetragenen Messungsgeschäfte dabei noch einige Verzögerung verursachen.
\par Das Schlusstheorem $b\equiv\frac12rr\, (\mathrm{mod}.p)$ hatte ich schon vor drei Jahren in den hiesigen gel. Anzeigen mit bekannt, und auf den merkwürdigen dabei noch zu lösenden Knoten aufmerksam gemacht; ich habe aber bisher nicht gehört, dass jemand einen Versuch dazu gemacht hätte. Vor einigen Tagen ist es mir nun mit der \textit{einen Hälfte} wirklich gelungen, und dieser Fund hat mir um so mehr Vergnügen gemacht, da er sich gar nicht auf Induction gründet \textemdash{} denn ich gestehe, dass ich gerade \textit{diesen} Zusammenhang nicht erwartet hätte \textemdash{} sondern a priori auf die Combination anderweitiger \textit{sehr verschlungener} und interessanter, schon 28 Jahr alter, aber noch gar nicht bekannt gemachter Untersuchungen,
